\scanpage{147}
\mspage{96}

L'idée naturelle, pour une réciproque : \unsure{prop 8|prop 7}, serait de prendre \inserted{pour $U'$} la somme amalgamée de $U$ et $V'$ sur $V$, ce qui \deleted{\illegible{2}} marchera pourvu qu'on sache que $\alpha'$ est bel et bien un iso. local tout comme $\alpha$. C'est bien le cas p. ex. si $V\unsure{-|\to}U$ est une imm. ouverte, i.e. $H$ un sous-préfaisceau du préfaisceau final.

Mais je crains que ce ne soit pas vrai en général -- ça a l'air un peu trop beau pour être vrai ! C'est en effet déjà faux dans le cas

\diagram{\begin{tikzcd}[ampersand replacement=\&]
V=E_G \arrow[r,"\alpha\ \mathrm{can}","\deleted{\unsure{g|f}}"'] \arrow[d,"f"'] \& B_G=U\\
V'=B_{G'} \&
\end{tikzcd}}

dans le cas p. ex. où $G=\{\pm1\}$ groupe à deux éléments, et où l'élément \inserted{$\alpha$} de $G'$ image par $f$ de l'\deleted{\unsure{flèches|identiques}} deux flèches non identiques de $E_G$ est $\ne1$. Si $H'$ est le sous-groupe invariant engendré par $\alpha$ dans $G'$, la \deleted{quotient} somme amalgamée s'identifie à $B_{G'/H'}$, or $B_{G'}\to B_{G'/H'}$ n'est \unsure{pas|$=$} un iso local (car $(B_{G'})_{/t_0}\to(B_{G'/H'})_{/t'_0}$ n'est pas bijectif sur les objets, s'identifiant à $G'\unsure{\nrightarrow|\to}G'/\unsure{H'|H}\ \unsure{???}$).

Mais d'autre part la condition de totale $\unsure{W|\mathcal W}$-asphéricité sur un objet $H$ de $U^{\unsure{\wedge|\sim}}$ \deleted{\unsure{par}}paraît déjà à tel point exigeante, que je ne vois mal \unsure{comm}
% archive sheet number: 147
